	
	\newpage
% =========================================================	
\section{Liste der Notationen}
% =========================================================

	\sbreak
	\begin{center}
		\begin{tabular}{cp{12cm}}
			\toprule
			Notation						&	Beschreibung\\
			\midrule
			$\reals$						&	Die reellen Zahlen. \\%The real numbers including the zero.\\
			$\naturals$						&	Die natürlichen Zahlen, Null enthaltend. \\
			$\complexes$					&	Menge der komplexen Zahlen. \\
			%$\vec{u}$						&	The vector $u$. \\
			$u_i$							&	$i$. Eintrag des Vektors $u$. \\
			$\bigO$							&	Big O Notation für Laufzeitklassen.\\
%			$|x|_n$							&	Die $l_n$ Norm von $x$. \\
			%$f(x)$							&	Eine Funktion, die von $x$ abhängig ist. \\
			$\abs{x}$						&	Absoluter Wert von $x$. \\
			$\norm{x}_2$					&	Normiertes $x$ ($\ell-2$ Norm). \\
			$f'(x)$							&	Die Funktion $f(x)$, einmal abgeleitet. \\
			$f^{(d)}(x)$					&	Die $d$. Ableitung von $f(x)$. \\
			% Matrix Stuff
			$A\trans$						&	Transponierte der Matrix $A$. \\
			$A\inv$							&	Das Inverse der Matrix $A$. \\
			$\diag(A)$						&	Vektor mit den Diagonaleinträgen der Matrix $A$. \\
			$\det(A)$						&	Die Determinante der Matrix $A$. \\
			% operationen
			$\rd(x)$						&	$x$, gerundet. \\
			$y \shouldEq x$					&	$y$ soll gleich $x$ sein. \\
			$y := x$						& 	$y$ ist definiert als $x$. \\
			$y \gg x$						&	$y$ viel größer als $x$. \\
			$y \ll x$						&	$y$ deutlich kleiner als $x$. \\
			$\exp(x)$						&	Andere Schreibweise für $e^x$. \\
			% komplexes zeug
			$\realPart(z)$					&	Realteil einer komplexen Zahl. \\
			$\imaginaryPart(z)$				&	Imaginärteil einer komplexen Zahl. \\
			$\compl{z}$						&	Das komplex konjugierte einer komplexen Zahl $z$. \\
			$\tilde{a}$						&	Eine durch Numerik verfälschte Variable $a$. \\
			\bottomrule
		\end{tabular}
	\end{center}
			
	\newpage
	\begin{center}
		\begin{tabular}{cp{12cm}}
			\toprule
			Notation						&	Beschreibung\\
			\midrule
			% Zahlendarstellungen
			%$\beta$							&	Zahlenbasis \\
			$n$								&	Anzahl Stückstellen oder Anzahl Bits (je nach Kapitel). \\
			% Kondition & Stabilität
			$\fabs$							&	Absoluter Fehler. \\
			$\frel$							&	Relativer Fehler. \\
			$\cond()$						&	Die relative Konditionszahl. \\
			$\maxError$						&	Maschinengenauigkeit (der maximal bei einer Rundung auftretende Fehler). \\
			op								&	Beliebige, exakte Operation. \\
			op$_\text{M}$					&	Eine beliebige, fehlerbehaftete Operation. \\
			% Orthogonal + QR
			$\unitM$						&	Die Einheitsmatrix. \\
			$G_\varphi$						&	Rotationsmatrix bezüglich des Winkels $\varphi$. \\
			% Interpolation
			$c_i$							&	Koeffizient mit Index $i$. \\
			$p(x)$							&	Interpolation (polynoms-) funktion. \\
			$L_j$							&	Lagrange Basis. \\
			$H_i$							&	Kubische Basispolynome. \\
			% Quadratur
			$H$								& 	Insgesamte Größe des zu interpolierenden / integrierenden Bereiches. \\
			$h$								&	Länge eines (oder aller) Bereiche zwischen zwei Stützstellen. \\
			$\QT$							&	Die Trapezregel oder Keplersche Fassregel. \\
			$\QTS$							&	Die Trapezsumme. \\
			$\QS$							&	Simpsonregel. \\
			$\QSS$							&	Simpsonsumme. \\
			% Fourier Transformation
			$\DFT$							&	Diskrete Fourier Transformation. \\
			$\IDFT$							&	Inverse diskrete Fourier Transformation. \\
			$\FFT$							&	Fast Fourier Transformation. \\
			$\BFO$							&	Butterfly Operator. \\
			% Fixpunkte
			$x*$							&	Fixpunkt. \\
			$\Phi(x)$						&	Iterationsvorschrift einer iterativen Methode. \\
			$\delta t$						&	Schrittweite eines Iterationsverfahrens. \\
			% Eigenwerte
			$\lambda(A)$					&	Eigenwerte der Matrix $A$. \\
			$\lambda_{\max}(A)$				&	Der größte Eigenwert der Matrix $A$. \\
			$\mu$							&	Shift einer Power Iteration. \\
			\bottomrule
		\end{tabular}
	\end{center}